\begin{proposition}[秩一割线方程的耦合单调性与导数--范数同一律]\label{prop:xi-rankone-secular-monotonicity-derivative-norm}
\leanverified{paper\_xi\_rankone\_secular\_monotonicity\_derivative\_norm}
取互异实数 \(t_0,\dots,t_q\) 与权重 \(w_i>0\)，并令
\[
D:=\mathrm{diag}(t_0,\dots,t_q),\qquad g:=(\sqrt{w_0},\dots,\sqrt{w_q})^{\mathsf T}.
\]
对任意 \(b>0\)，考虑秩一扰动
\[
A(b):=D+b\,gg^{\mathsf T}.
\]
则 \(A(b)\) 的谱由割线方程
\begin{equation}\label{eq:xi-rankone-secular-general}
\Phi(\lambda):=\sum_{i=0}^{q}\frac{w_i}{\lambda-t_i}=\frac1b,
\qquad
\lambda\notin\{t_0,\dots,t_q\},
\end{equation}
给出，且任意根分支 \(\lambda=\lambda(b)\) 在其定义域内满足严格单调与导数闭式
\begin{equation}\label{eq:xi-rankone-secular-derivative}
\frac{\mathrm d\lambda}{\mathrm db}
=\frac{1}{b^2\displaystyle\sum_{i=0}^{q}\frac{w_i}{(\lambda-t_i)^2}}
\;>\;0.
\end{equation}
并且与 \(\lambda(b)\) 对应的特征向量可取为
\[
v(b):=(\lambda I-D)^{-1}g,\qquad
v_i(b)=\frac{\sqrt{w_i}}{\lambda(b)-t_i},
\]
从而耦合量与范数满足恒等式
\[
g^{\mathsf T}v(b)=\frac1b,
\qquad
\|v(b)\|_2^2=\sum_{i=0}^{q}\frac{w_i}{(\lambda(b)-t_i)^2}=\frac{1}{b^2\,\lambda'(b)}.
\]
\end{proposition}
\begin{proof}
由矩阵行列式引理，
\[
\det(\lambda I-A(b))
=\det(\lambda I-D)\Bigl(1-b\,g^{\mathsf T}(\lambda I-D)^{-1}g\Bigr),
\]
其中
\(
g^{\mathsf T}(\lambda I-D)^{-1}g=\sum_i\frac{w_i}{\lambda-t_i}=\Phi(\lambda)
\)。
故 \(\lambda\) 为特征值当且仅当 \eqref{eq:xi-rankone-secular-general} 成立；并且 \(\Phi'(\lambda)=-\sum_i\frac{w_i}{(\lambda-t_i)^2}<0\)，根均为单根。
对 \(\Phi(\lambda(b))=1/b\) 作隐式微分即得 \eqref{eq:xi-rankone-secular-derivative}，从而 \(\lambda'(b)>0\)。
\par
取 \(v(b)=(\lambda I-D)^{-1}g\)，则
\[
(A(b)-\lambda I)v
=(D-\lambda I)v+b\,g(g^{\mathsf T}v)
=-g+b\,g\,\Phi(\lambda)
=0,
\]
其中最后一步使用割线方程；故 \(v(b)\) 为特征向量。
其坐标式与 \(g^{\mathsf T}v=1/b\) 由定义直接给出；
范数恒等式由坐标式与 \eqref{eq:xi-rankone-secular-derivative} 联立得到。
\end{proof}

\begin{theorem}[秩一耦合的谱纤维分裂：小耦合与大耦合极限]\label{thm:xi-rankone-secular-small-large-coupling}
沿用命题 \ref{prop:xi-rankone-secular-monotonicity-derivative-norm} 的记号。
令极点集合 \(t_i\) 两两不同，并记
\[
G(\lambda):=\prod_{i=0}^{q}(\lambda-t_i),
\qquad
N(\lambda):=\sum_{i=0}^{q} w_i\!\!\prod_{\substack{0\le j\le q\\ j\ne i}}\!\!(\lambda-t_j),
\qquad
\Phi(\lambda)=\frac{N(\lambda)}{G(\lambda)}.
\]
则对 \(b>0\) 的全部特征值（即 \(\Phi(\lambda)=1/b\) 的全部根）有：
\begin{enumerate}
  \item 当 \(b\to 0^+\) 时，对每个 \(i\in\{0,\dots,q\}\) 存在唯一根分支 \(\lambda_i(b)\) 满足
  \[
  \lambda_i(b)=t_i+bw_i+O(b^2),
  \qquad
  \lambda_i(b)\to t_i.
  \]
  \item 当 \(b\to+\infty\) 时，存在唯一根分支 \(\lambda_\infty(b)\to+\infty\) 且
  \[
  \lambda_\infty(b)=b\sum_{i=0}^{q}w_i+\frac{\sum_{i=0}^{q}w_it_i}{\sum_{i=0}^{q}w_i}+O(b^{-1}).
  \]
  其余 \(q\) 条根分支在 \(b\to+\infty\) 时分别收敛到多项式 \(N(\lambda)\) 的 \(q\) 个实零点；这些零点均为单根并严格夹在相邻极点之间。
\end{enumerate}
\end{theorem}
\begin{proof}
\par\noindent\textbf{(i) 小耦合.}
固定 \(i\)。在 \(\lambda=t_i\) 的穿孔邻域内，
\[
\Phi(\lambda)=\frac{w_i}{\lambda-t_i}+H_i(\lambda),
\qquad
H_i(\lambda):=\sum_{j\ne i}\frac{w_j}{\lambda-t_j},
\]
其中 \(H_i\) 在 \(\lambda=t_i\) 解析。
令 \(\lambda=t_i+\delta\)，方程 \(\Phi(\lambda)=1/b\) 等价于
\(
w_i/\delta+H_i(t_i)+O(\delta)=1/b
\)。
当 \(b\to 0^+\) 时右端趋于 \(+\infty\)，迫使 \(\delta\to 0\)，并由上式解得
\(\delta=bw_i+O(b^2)\)。
由于 \(\Phi'(\lambda)<0\)，在足够小的邻域内根唯一，从而得到唯一根分支。
\par
\noindent\textbf{(ii) 大耦合.}
当 \(|\lambda|\to\infty\)，有渐近展开
\[
\Phi(\lambda)=\sum_{i=0}^{q}\frac{w_i}{\lambda-t_i}
=\frac{W}{\lambda}+\frac{M_1}{\lambda^2}+O(\lambda^{-3}),
\qquad
W:=\sum_{i=0}^{q}w_i,\quad M_1:=\sum_{i=0}^{q}w_it_i.
\]
令 \(\Phi(\lambda)=1/b\) 并令 \(b\to+\infty\) 得到唯一逃逸根满足 \(\lambda\sim bW\)；
进一步迭代代回即得所示渐近展开。
\par
其余根的收敛：注意 \(N(\lambda)=0\) 当且仅当 \(\Phi(\lambda)=0\)。
在每个相邻极点之间的开区间上，\(\Phi\) 连续且严格单调（\(\Phi'<0\)），并在端点趋于 \(\pm\infty\)，因此 \(N\) 在每个此类区间内恰有一个实零点且为单根。
同理，\(\Phi(\lambda)=1/b\) 在每个区间内亦恰有一个实根；令 \(b\to+\infty\) 使右端趋于 \(0\)，由单调性与连续性得该根收敛到区间内的零点。
\end{proof}

\endinput
